% Proof supplement for admissible tag reassignment.
% Root manuscript should supply bibliography entries for the inherited
% RestrictedImplication proof and the radix-family calibration.
\section{Formula-tag reassignment without constructor collisions}
\label{sec:tag-reassignment}

Changing the residue used by equality or universal quantification is a
mathematical change to the numerical calculus, not a state renaming. The
following conditions make that change sound. They also explain why changing
the multiplier to three while retaining four constructor tags is invalid.

\subsection{Admissible tags}

Fix $k\ge4$ and two distinct numbers $u,e$ with $2\le u,e<k$.
For Cantor pairing $p(a,b)=((a+b)(a+b+1))/2+b$, use
\begin{align*}
 I(A,B)&=kp(A,B), & M(a,b)&=kp(a,b)+1,\\
 U(a,B)&=kp(B,a)+u, & E(a,b)&=kp(a,b)+e.
\end{align*}
Code zero is truth and code one is the designated contradiction $M(0,0)$.
Negation, conjunction and existential quantification use the same
implicational definitions as the baseline calculus. Retain the six logical
heads, the dominated implicational head, the separate diagonal head and the
guarded separation gate, replacing their constructors by the displayed ones.
The guard remains $c\ne a,b$.

The relevant instances include $(k,u,e)=(4,3,2)$ and $(5,2,4)$. The first
swaps the two constructor residues. The second leaves residue three unused
and moves equality to residue four. In each case all four constructor tags
are distinct and the implication and contradiction conventions remain fixed.

Every component of a nonzero code is strictly smaller than the code. Indeed,
if $n=kp(a,b)+r>0$ then $a,b\le p(a,b)<n$; when $r=0$, positivity implies
$p(a,b)>0$, and when $p(a,b)=0$, both components are zero and $r>0$.
Consequently the usual course-of-values coding recursions are primitive
recursive. The identity $I(0,0)=0$ is still exceptional and must not be
treated as disjoint syntax for truth and implication.

\subsection{An explicit hereditary embedding}

Let $\tau(0)=0$, $\tau(1)=1$, $\tau(2)=u$, $\tau(3)=e$. Define $F(0)=0$;
for $n>0$, uniquely write $n=4p(a,b)+r$ and put
\[
 F(n)=kp(F(a),F(b))+\tau(r).
\]
The function is primitive recursive: compute its values through $n$ in an
accumulated finite table, looking up only smaller arguments. It fixes zero
and one and maps positive inputs to positive outputs. Define a left inverse
$J$ by the analogous recursion on target codes: use $\tau^{-1}(r)$ for
the four used residues, and return zero at an unused residue. Induction on
the old code gives $J(F(n))=n$, so $F$ is injective.

Both formula codes and numerical variable labels are transformed by this
same map. Directly from its defining equations,
\begin{align*}
 F(I_4(A,B))&=I(F(A),F(B)),\\
 F(M_4(a,b))&=M(F(a),F(b)),\\
 F(U_4(a,B))&=U(F(a),F(B)),\\
 F(E_4(a,b))&=E(F(a),F(b)).
\end{align*}
The first equation includes the exceptional zero case. Injectivity preserves
guards, binding and repeated labels. Most importantly, it preserves the
diagonal head exactly:
\[
 F\bigl(I_4(B,U_4(B,B))\bigr)
   =I\bigl(F(B),U(F(B),F(B))\bigr).
\]
Recoding formula constructors while leaving their numerical labels fixed
would not establish this equation.

The embedding need not preserve numerical order. In the swap case,
$F(2)=3$ and $F(3)=2$. Therefore an old restricted propositional leaf cannot
be translated as one new restricted leaf merely by applying $F$ to its
parameters.

\subsection{The dominated implicational head is no obstruction}

Write
\[
 \rho(D,A,H,d)=((H\to d)\to D)\to((D\to H)\to(A\to H)).
\]
The machine emits this head when $d\le D$. The existing finite proof
\texttt{rho8All} in \path{M6/Logic/RestrictedImplication.lean}
\cite{restricted-implication} derives the
unrestricted head from restricted instances, using only modus ponens,
reflexivity of the side relation, and the two implications
\[
 d\le A\Longrightarrow d\le I(A,B),\qquad
 d\le B\Longrightarrow d\le I(A,B).
\]
These hold here because $A,B\le p(A,B)\le kp(A,B)$. The proof is generic
in the implication operation and contains no quantifier, membership or
equality constructor. Thus changing $u,e$ does not change its validity.

Translate every other leaf and inference by $F$, and replace each translated
propositional leaf by that unrestricted finite template. This gives a
primitive-recursive translation of baseline-calculus derivations into the
new calculus, preserving contradiction since $F(1)=1$. The domination
conditions of the template's own leaves are proved directly in the target
arithmetic. No monotonicity or subtraction identity for $F$ is used.

For $(k,u,e)=(4,3,2)$, the same hereditary map is an involution. Its defining
tag permutation is its own inverse, and induction using the component
equations proves $F(F(n))=n$. Applying the same proof translation in the
opposite direction gives refutation translations between the two calculi.
For larger $k$, $J$ is only a left inverse on the image; it is not claimed
to translate arbitrary target proofs back, since unused tags and possible
label collisions prevent that inference.

\subsection{Soundness of all target numerical codes}

To include unused residues, define a total syntactic interpretation $T$ in
arithmetic. Send zero and every unused-residue code to closed truth. For a
positive code $n=kp(a,b)+r$, the used residues are interpreted by
\[
\begin{array}{c|l}
r& T(n)\\\hline
0&T(a)\to T(b)\\
1&\operatorname{Bit}(x_a,x_b)\\
u&\forall x_b\,T(a)\\
e&x_a=x_b.
\end{array}
\]
Auxiliary arithmetic bound variables are kept separate from the source
labels by de Bruijn coding or capture-avoiding renaming. Every source label
occurring in $T(B)$ is strictly below $B$, by the component inequality.
An unused residue introduces no source label. Hence the diagonal formula
$B\to U(B,B)$ has a uniformly generated PA proof for every numeric $B$,
including junk codes.

The other axiom suppliers are unchanged in mathematical content. Equality
congruence and its existential witness are ordinary first-order equality
theorems; quantifier distribution, commutation and witness elimination have
their usual proof templates. The separation supplier uses finite binary
comprehension for the syntactically translated formula $T(c)$, with the
distinct-label and coincident-label cases handled exactly as in the
baseline calibration \cite{radix-parent}. The guard and the freshness bound are the only
properties of the numerical tags used in that argument.

The implication translation is a syntactic identity except at $I(0,0)=0$,
where a fixed proof of
$(0=0)\leftrightarrow((0=0)\to(0=0))$ supplies the required equivalence.
The interpretation of code one is $\operatorname{Bit}(x_0,x_0)$, which PA
refutes. Recursing over a target derivation and inserting the appropriate
proof supplier therefore gives a primitive-recursive translation of target
refutations into PA refutations, with the same paper-proof scope as the
inherited calibration. No uniform arithmetic truth predicate is used.

\subsection{The scheduler and stack offsets must be recompiled}

The stack continues to use ordinary Cantor pairing, not the formula
constructors. An AX command pushes the same ten entries. Its guarded top
and the eight leading logical entries have root implication tag zero.
The generalization entry $U(d,h)$ has residue $u\ne0$ and value at least
two. These are precisely the residue and size facts needed by the safe
SHIFT windows, PUSH, POP and DISCARD words. Replacing the old universal
tag two by any allowed $u$ therefore preserves those proofs.

For a target restricted leaf with parameters $d\le D$, the compiler uses
the actual target offset $D-d$, so $d+(D-d)=D$. The fixed unrestricted-rho
expansion just described supplies leaves meeting that restriction. The MP
compiler similarly calculates offsets in the target numerical encoding.
The formula embedding preserves neither ordinary addition nor Cantor-stack
coding; identities such as $F(d+x)=F(d)+F(x)$ are neither true in general
nor needed. One compiles a new target command word from the translated
derivation, and then invokes the scheduler's exhaustive word search.

Finally, the source generator must be checked to implement these constructors
at \emph{every} occurrence. Altering only some equality tags would not be
a tagged-calculus implementation. A free-constructor emitter check, exact
primitive source-map binding, positive full-frame physical contracts and
whole-tape blank initialization are distinct obligations. Successful tag
algebra does not replace any of those machine checks.

The new tag argument can therefore be combined with the inherited written
PA calibration, or with a future completed iPRA calibration after its stated
syntactic and physical obligations are discharged. It is not, by itself,
an end-to-end Lean proof or a completion of the baseline's open iPRA bridge.
